% =========================================================
\section{Interlude: Memory API}
% =========================================================

\begin{frame}{Memory API}
    \begin{block}{The Crux: How to Allocate and Manage Memory}
        In UNIX/C programs, understanding how to allocate and manage
        memory is critical for building robust and reliable software.
    \end{block}

    \begin{itemize}
        \item What interfaces are commonly used?
        \item What mistakes should be avoided?
    \end{itemize}
\end{frame}


\begin{frame}[fragile]{Types of Memory}
    \begin{columns}
        \begin{column}{0.48\textwidth}
            \textbf{Stack Memory}
            \begin{itemize}
                \item Managed implicitly by the compiler.
                \item Also called \emph{automatic memory}.
                \item Allocated when a function is called.
                \item Deallocated when the function returns.
            \end{itemize}

\begin{footnotesize}
\begin{verbatim}
void func() {
    int x;
    ...
}
\end{verbatim}
\end{footnotesize}
        \end{column}

        \begin{column}{0.48\textwidth}
            \textbf{Heap Memory}
            \begin{itemize}
                \item Explicitly managed by the programmer.
                \item Used when memory must outlive a function call.
                \item Allocation and deallocation require explicit calls.
            \end{itemize}

\begin{footnotesize}
\begin{verbatim}
void func() {
    int *x =
      (int *) malloc(sizeof(int));
    ...
}
\end{verbatim}
\end{footnotesize}
        \end{column}
    \end{columns}
\end{frame}


\begin{frame}[fragile]{The \texttt{malloc()} Call}
    \begin{block}{Interface}
\begin{verbatim}
#include <stdlib.h>
void *malloc(size_t size);
\end{verbatim}
    \end{block}

    \begin{itemize}
        \item Requests \texttt{size} bytes of heap memory.
        \item Returns a pointer to the allocated region on success.
        \item Returns \texttt{NULL} on failure.
        \item \texttt{sizeof()} is commonly used to determine
              the required allocation size.
    \end{itemize}

\begin{verbatim}
double *d = (double *) malloc(sizeof(double));
int *x = malloc(10 * sizeof(int));
\end{verbatim}
\end{frame}


\begin{frame}[fragile]{Using \texttt{sizeof()}}
    \begin{itemize}
        \item \texttt{sizeof()} is an operator, generally evaluated
              at compile time.
        \item Be careful when applying it to pointers.
    \end{itemize}

\begin{verbatim}
int *x = malloc(10 * sizeof(int));
printf("%d\n", sizeof(x));
\end{verbatim}

    \begin{itemize}
        \item \texttt{sizeof(x)} returns the size of the
              \emph{pointer}, not the dynamically allocated region.
    \end{itemize}

\begin{verbatim}
int x[10];
printf("%d\n", sizeof(x));
\end{verbatim}

    \begin{itemize}
        \item Here, the compiler has enough information to determine
              the size of the entire array.
    \end{itemize}
\end{frame}


\begin{frame}{Tip: When in Doubt, Try It Out}
    \begin{block}{Tip}
        If you are unsure how a routine or operator behaves,
        write some code and test it.
    \end{block}

    \begin{itemize}
        \item Manual pages and documentation are useful.
        \item Testing confirms how something behaves in practice.
        \item The chapter uses this approach to verify the behavior
              of \texttt{sizeof()}.
    \end{itemize}
\end{frame}


\begin{frame}[fragile]{Strings and \texttt{malloc()}}
    \begin{itemize}
        \item String allocation requires room for the
              end-of-string character.
        \item The suggested idiom is:
    \end{itemize}

\begin{verbatim}
malloc(strlen(s) + 1)
\end{verbatim}

    \begin{itemize}
        \item \texttt{strlen()} determines the string length.
        \item The additional byte stores the terminating character.
        \item Using \texttt{sizeof()} for this purpose may lead
              to errors.
        \item The cast of the \texttt{void *} returned by
              \texttt{malloc()} is not required for correctness.
    \end{itemize}
\end{frame}


\begin{frame}[fragile]{The \texttt{free()} Call}
    \begin{itemize}
        \item Allocating memory is relatively easy;
              knowing when and how to free it is harder.
        \item Heap memory that is no longer needed should be
              released with \texttt{free()}.
    \end{itemize}

\begin{verbatim}
int *x = malloc(10 * sizeof(int));
...
free(x);
\end{verbatim}

    \begin{itemize}
        \item \texttt{free()} receives a pointer previously
              returned by \texttt{malloc()}.
        \item The allocation size is tracked by the
              memory-allocation library.
    \end{itemize}
\end{frame}


\begin{frame}{Common Memory Errors}
    \begin{itemize}
        \item Forgetting to allocate memory.
        \item Not allocating enough memory (\emph{buffer overflow}).
        \item Forgetting to initialize allocated memory.
        \item Forgetting to free memory (\emph{memory leak}).
        \item Freeing memory before use is complete
              (\emph{dangling pointer}).
        \item Freeing memory repeatedly (\emph{double free}).
        \item Calling \texttt{free()} with an invalid pointer.
    \end{itemize}

    \begin{alertblock}{Important}
        Compiling successfully is necessary, but far from sufficient,
        for a correct C program.
    \end{alertblock}
\end{frame}


\begin{frame}[fragile]{Forgetting to Allocate Memory}
\begin{verbatim}
char *src = "hello";
char *dst;                 // unallocated
strcpy(dst, src);          // likely segfault
\end{verbatim}

    \begin{itemize}
        \item Some routines expect the destination memory to have
              already been allocated.
    \end{itemize}

\begin{verbatim}
char *src = "hello";
char *dst =
    (char *) malloc(strlen(src) + 1);
strcpy(dst, src);
\end{verbatim}

    \begin{itemize}
        \item Alternatively, \texttt{strdup()} can be used.
    \end{itemize}
\end{frame}


\begin{frame}{Tip: It Compiled or Ran $\neq$ It Is Correct}
    \begin{block}{Tip}
        A program that compiles, or even runs correctly many times,
        is not necessarily correct.
    \end{block}

    \begin{itemize}
        \item Memory errors may remain unnoticed during some executions.
        \item Small changes can expose previously hidden bugs.
        \item Buffer overflows are an important example:
              they may appear harmless, crash the program,
              or create serious security problems.
    \end{itemize}
\end{frame}


\begin{frame}{Memory Leaks and Dangling Pointers}
    \begin{columns}
        \begin{column}{0.48\textwidth}
            \textbf{Memory Leak}
            \begin{itemize}
                \item Allocated memory is never freed.
                \item Particularly problematic for long-running
                      applications and systems.
                \item Eventually, the application may run out
                      of memory.
            \end{itemize}
        \end{column}

        \begin{column}{0.48\textwidth}
            \textbf{Dangling Pointer}
            \begin{itemize}
                \item Memory is freed before the program
                      has finished using it.
                \item Subsequent use may crash the program
                      or overwrite valid memory.
            \end{itemize}
        \end{column}
    \end{columns}
\end{frame}


\begin{frame}{Aside: Why No Memory Is Leaked Once a Process Exits}
    \begin{itemize}
        \item There are two levels of memory management:
        \begin{enumerate}
            \item the OS manages memory assigned to processes;
            \item each process manages allocations within its own heap.
        \end{enumerate}

        \item A process may fail to call \texttt{free()} and thus
              leak memory within its heap.

        \item When the process exits, however, the OS reclaims
              all of its memory:
        \begin{itemize}
            \item code;
            \item stack;
            \item heap.
        \end{itemize}

        \item Thus, leaks are especially problematic in
              \textbf{long-running programs} and the \textbf{OS itself}.
    \end{itemize}
\end{frame}


\begin{frame}{More Memory Errors}
    \begin{description}
        \item[Double free]
        Calling \texttt{free()} more than once on the same memory;
        the result is undefined.

        \item[Invalid free]
        Passing to \texttt{free()} a value other than an appropriate
        pointer previously returned by \texttt{malloc()}.
    \end{description}

    \begin{block}{Finding Memory Problems}
        Tools such as \texttt{purify} and \texttt{valgrind}
        can help locate memory-related errors.
    \end{block}
\end{frame}


\begin{frame}{Underlying OS Support}
    \begin{itemize}
        \item \texttt{malloc()} and \texttt{free()} are
              \textbf{library calls}, not system calls.
        \item The allocation library manages memory within the
              process's virtual address space.
        \item It relies on system calls to obtain or release memory.
    \end{itemize}

    \begin{description}
        \item[\texttt{brk}/\texttt{sbrk}]
        Change the end of the heap; applications should not call
        them directly.

        \item[\texttt{mmap()}]
        Can create an anonymous memory region that can also be
        managed like heap memory.
    \end{description}
\end{frame}


\begin{frame}{Other Memory Allocation Calls}
    \begin{description}
        \item[\texttt{calloc()}]
        Allocates memory and initializes it to zero.

        \item[\texttt{realloc()}]
        Creates a new, larger memory region, copies the old region
        into it, and returns a pointer to the new region.
    \end{description}

    \begin{block}{Key Idea}
        These routines complement \texttt{malloc()} and
        \texttt{free()} for common dynamic-memory management needs.
    \end{block}
\end{frame}


\begin{frame}{Summary}
    \begin{itemize}
        \item C programs use both \textbf{stack} and \textbf{heap} memory.
        \item Heap memory is explicitly managed by the programmer.
        \item \texttt{malloc()} allocates memory and
              \texttt{free()} releases it.
        \item Incorrect memory management can cause:
              buffer overflows, uninitialized reads, memory leaks,
              dangling pointers, double frees, and invalid frees.
        \item \texttt{malloc()} and \texttt{free()} are library calls
              built on top of OS memory-management mechanisms.
        \item Tools such as \texttt{valgrind} help identify
              memory-related problems.
    \end{itemize}
\end{frame}